\documentclass[aspectratio=169,11pt]{beamer}

\usetheme{Madrid}
\usecolortheme{default}
\usefonttheme{professionalfonts}
\setbeamertemplate{navigation symbols}{}
\setbeamertemplate{footline}[frame number]

\usepackage{amsmath, amssymb, booktabs, graphicx}

\title{Identity, Social Hierarchy, and Labor-Market Institutions}
\subtitle{Labor Markets and Political/Cultural Institutions --- Week 6}
\author{Teaching deck draft}
\date{}

\begin{document}

\begin{frame}
  \titlepage
\end{frame}

\begin{frame}{Central question}
\begin{itemize}
    \item How do institutionalized social categories structure access to jobs, networks, authority, protection, and mobility?
    \item When do caste, race, ethnicity, gender, religion, migrant status, or related categories become labor-market institutions?
    \item Which evidence identifies employer beliefs, network closure, public rules, contact, or misallocation?
\end{itemize}
\end{frame}

\begin{frame}{Bridge from Week 5}
\begin{itemize}
    \item Week 5: norms of conduct and acceptable labor-market behavior.
    \item Week 6: durable category-based hierarchy and institutionalized access.
    \item Conduct norm question: what action is rewarded or punished?
    \item Hierarchy question: how does category membership structure access, authority, protection, and mobility?
\end{itemize}
\end{frame}

\begin{frame}{Week 5 versus Week 6}
\small
\begin{tabular}{p{0.30\linewidth} p{0.30\linewidth} p{0.30\linewidth}}
\toprule
Object & Week 5 & Week 6 \\
\midrule
Main mechanism & conduct and role expectations & category-based access and exclusion \\
Examples & family roles, competition, reputation, wage acceptance & caste, race, migrant status, religion, gender hierarchy \\
Labor margin & labor supply, search, entry, acceptable jobs & callbacks, referrals, credentials, protection, occupations \\
Diagnostic & what conduct is costly? & which gate is category-ranked? \\
\bottomrule
\end{tabular}
\end{frame}

\begin{frame}{When categories become institutions}
\begin{itemize}
    \item A category becomes institutional when it repeatedly changes labor-market gates.
    \item Gates include hiring, promotion, credential recognition, network access, occupation entry, public jobs, and legal protection.
    \item The category matters because organizations and rules make it predictive of access.
    \item Start from the labor outcome, then name the institutional channel.
\end{itemize}
\end{frame}

\begin{frame}{A compact access map}
\[
P(\text{access}_{ij}) = f(x_i, q_i, B_j, N_i, L_i, H_i)
\]
\begin{itemize}
    \item $x_i$: observed human capital.
    \item $q_i$: job-relevant quality that may be imperfectly observed.
    \item $B_j$: employer beliefs or screening rules.
    \item $N_i$: network access; $L_i$: legal or public-rule status.
    \item $H_i$: category-based hierarchy.
\end{itemize}
\end{frame}

\begin{frame}{Employer beliefs and differential treatment}
\begin{itemize}
    \item Resume names, accents, addresses, schools, and credentials become screening signals.
    \item Bertrand--Mullainathan: randomized names on resumes identify callback differential treatment.
    \item Oreopoulos: skilled immigrants, foreign credentials, and employer interpretation of background signals.
    \item Strongest claim: first-stage access changes when the category signal changes.
\end{itemize}
\end{frame}

\begin{frame}{What audit designs identify}
\small
\begin{tabular}{p{0.30\linewidth} p{0.29\linewidth} p{0.29\linewidth}}
\toprule
Design & Identifies best & Does not identify alone \\
\midrule
Resume audit & differential treatment at a hiring gate & wages, promotions, lifecycle effects \\
Credential experiment & value assigned to visible credentials & true productivity, licensing, long-run match quality \\
Name or accent signal & screening response to category signal & taste versus statistical discrimination without more evidence \\
\bottomrule
\end{tabular}
\end{frame}

\begin{frame}{Networks, referrals, and occupation persistence}
\begin{itemize}
    \item Networks transmit information, insure risk, screen workers, and enforce obligations.
    \item They reproduce hierarchy when valuable referrals are category-closed.
    \item Munshi--Rosenzweig: caste networks, schooling, occupations, and mobility.
    \item Key diagnostic: is the network solving information problems, restricting mobility, or both?
\end{itemize}
\end{frame}

\begin{frame}{Credential recognition and occupation boundaries}
\begin{itemize}
    \item Skills can have high productive value but low market value if credentials are discounted.
    \item Migrant status often combines name screening, local-experience requirements, licensing, and legal exposure.
    \item Occupation boundaries restrict entry into high-return tracks and authority positions.
    \item The labor object is underemployment, downgrading, promotion barriers, and constrained matching.
\end{itemize}
\end{frame}

\begin{frame}{Talent misallocation}
\[
Y = Y^*(1 - \mu_H)
\]
\begin{itemize}
    \item $Y^*$: output under a less distorted allocation of talent.
    \item $\mu_H$: output share lost because hierarchy restricts occupational access.
    \item Hsieh--Hurst--Jones--Klenow: barriers by race and gender reduce aggregate talent allocation.
    \item Audit evidence does not mechanically scale to $\mu_H$; the bridge needs equilibrium assumptions.
\end{itemize}
\end{frame}

\begin{frame}{Public rules and state action}
\begin{itemize}
    \item Hierarchy can be produced by the state, not only by private employers.
    \item Public employment, legal status, licensing, enforcement, schools, housing, and policing shape labor access.
    \item Aneja--Xu: federal employment segregation under Woodrow Wilson.
    \item Historical and public-rule shocks identify effects of institutional reorganization on labor outcomes.
\end{itemize}
\end{frame}

\begin{frame}{Local governance and spatial sorting}
\begin{itemize}
    \item Local implementation determines who gets protection, whose complaints are credible, and which firms comply.
    \item Residential sorting changes commuting costs, job information, school access, and exposure to employers.
    \item Migrant settlement and ethnic enclaves can lower job-search costs while segmenting occupations.
    \item Place is an institutional margin when it structures access and protection.
\end{itemize}
\end{frame}

\begin{frame}{Contact and integration}
\begin{itemize}
    \item Contact is an institutional design problem, not just proximity.
    \item Lowe: collaborative and adversarial caste contact have different effects.
    \item Workplace teams, training programs, and public works projects can reduce or reinforce hierarchy.
    \item Contact experiments identify the effect of a designed interaction, not all labor-market barriers.
\end{itemize}
\end{frame}

\begin{frame}{Identification map}
\small
\begin{tabular}{p{0.31\linewidth} p{0.29\linewidth} p{0.28\linewidth}}
\toprule
Design & Identifies best & Main caution \\
\midrule
Audit or resume experiment & differential treatment at a gate & partial labor margin \\
Network or migration design & referral access and persistence & efficiency versus exclusion \\
Public-rule shock & state-created access changes & spillovers and persistence \\
Contact experiment & effects of interaction design & external validity \\
Misallocation framework & aggregate cost of barriers & model dependence \\
\bottomrule
\end{tabular}
\end{frame}

\begin{frame}{Micro, meso, macro}
\begin{itemize}
    \item Micro: callbacks, interviews, credential screening, individual treatment.
    \item Meso: networks, local governance, occupation boundaries, spatial sorting.
    \item Macro: aggregate talent allocation, output losses, stratification across groups.
    \item Frontier work links levels without treating one estimate as the whole mechanism.
\end{itemize}
\end{frame}

\begin{frame}{Week 6 lab}
\begin{itemize}
    \item Reproduce: synthetic resume-audit exercise inspired by Bertrand--Mullainathan.
    \item Diagnose: classify barriers as belief, network, credential, legal, contact, spatial, or macro.
    \item Transfer: map designs to Oreopoulos, Munshi--Rosenzweig, Lowe, Aneja--Xu, Hsieh et al., and Darity.
    \item Core output: gate, institutional channel, evidence level, and Week 5 versus Week 6 boundary.
\end{itemize}
\end{frame}

\begin{frame}{Frontier questions}
\begin{itemize}
    \item How do micro hiring barriers become macro allocation losses?
    \item Which institutions weaken hierarchy rather than moving it to another gate?
    \item How do legal status, credential recognition, and labor demand interact for migrants?
    \item When do public employers mitigate hierarchy, and when do they entrench it?
    \item Which forms of contact reduce hierarchy durably?
\end{itemize}
\end{frame}

\begin{frame}{Bridge to Week 7}
\begin{itemize}
    \item Week 6 studies hierarchy as a current allocation system.
    \item Week 7 asks why past institutions persist after formal rules change.
    \item Persistence can run through schooling, occupations, assets, places, networks, and public employment histories.
    \item The next question: when does institutional origin become durable labor-market inequality?
\end{itemize}
\end{frame}

\begin{frame}{Takeaways}
\begin{itemize}
    \item Category hierarchy becomes institutional when it governs access to labor-market gates.
    \item Main channels: beliefs, networks, credentials, occupations, public rules, local governance, place, and contact.
    \item Designs identify different objects; name the gate before interpreting the gap.
    \item The frontier connects micro barriers to meso persistence and macro talent allocation.
\end{itemize}
\end{frame}

\end{document}
